\documentclass{article}
\usepackage[T1]{fontenc}
\usepackage{lmodern}
\usepackage{graphicx}
\usepackage{amsmath,amssymb}
\usepackage[a4paper,margin=2cm]{geometry}
\usepackage[absolute,overlay]{textpos}
\setlength{\TPHorizModule}{1mm}
\setlength{\TPVertModule}{1mm}
\usepackage[indonesian]{babel}
\usepackage{hyperref}
\setlength{\emergencystretch}{2em}
% unit: Halaman 3

\begin{document}

% === Halaman 3 (python/native) ===

3

Blok plainteks berukuran n bit:

 P = (p1, p2, …, pn), pi  \{0, 1\} 

Enkripsi, E

p1

p2

...

pn

c1

c2

...

cn

Kunci K

Blok plainteks P

Blok cipherteks C

Dekripsi, D

c1

c2

...

cn

p1

p2

...

pn

Kunci K

Blok plainteks P

Blok cipherteks C

Blok cipherteks berukuran n bit:

 C = (c1, c2, …, cn), ci  \{0, 1\}

\end{document}
